% GENERATED FILE. Edit canonical lessons, then run npm run generate:books.
% canonical-source-hash: fnv1a64:2d7a9cda793cd41b
% canonical-lessons: MR-C122-sadhan, MR-C122-malak, MR-C122-takad, MR-C122-hava, MR-C122-iccha

\chapter{Means, Owner, Strength, Wanted, Wish}
\label{ch:mr-means-owner-strength-wanted-wish}
\hlchaptermodality{122}

\begin{chapteropening}
\textbf{By the end of this chapter:} \emph{I can say a means, an owner, strength, wanted and a wish.}

Complete the last lesson of chapter 122: I can say a means, an owner, strength, wanted and a wish.
\end{chapteropening}

\section[\texorpdfstring{\mr{साधन} (sādhan)}{साधन (sādhan)}]{\mr{साधन} (sādhan) — a means, a tool}

\label{lesson:MR-C122-sadhan}



\begin{quote}
Before the new one: say the Marathi for ability, then the Marathi for a right.
\end{quote}

\subsection*{You'll want to know: \mr{साधन}}
\textbf{\mr{साधन}} — \emph{sādhan} — \textquotedblleft{}a means, a tool\textquotedblright{}.

\begin{cousinweb}
\textbf{\mr{साधन}} (\emph{sādhan}) is Sanskrit \emph{sādhana}, \textquotedblleft{}furthering; completion\textquotedblright{}, taken over whole as a learned loan, a tatsama (CDIAL 13334). It is built on \emph{sādh}, \textquotedblleft{}to accomplish\textquotedblright{}.
\end{cousinweb}

\begin{grammarlens}[title={What you've built}]
One more word for this chapter.
\end{grammarlens}

\subsection*{Your turn}
\begin{itemize}
\raggedright
  \item \emph{Say it:} \emph{sādhan}
  \item \emph{Say it:} \emph{sādhan}, once more
  \item \emph{Say it:} say \emph{sādhan}
  \item \emph{Recall:} say the Marathi for ability, then the Marathi for a right, then say \emph{sādhan} again
\end{itemize}

\begin{tcolorbox}[breakable,colback=teal!4,colframe=teal!35!black,title={Before you move on}]
What does \mr{साधन} mean? (\textquotedblleft{}A means, a tool\textquotedblright{}.) Say it once more.
\end{tcolorbox}
\section[\texorpdfstring{\mr{मालक} (mālak)}{मालक (mālak)}]{\mr{मालक} (mālak) — an owner}

\label{lesson:MR-C122-malak}



\begin{quote}
Before the new one: say the Marathi for a right, then the Marathi for a means.
\end{quote}

\subsection*{You'll want to know: \mr{मालक}}
\textbf{\mr{मालक}} — \emph{mālak} — \textquotedblleft{}an owner\textquotedblright{}.

\begin{cousinweb}
\textbf{\mr{मालक}} (\emph{mālak}) is a loan from Arabic \emph{mālik}, \textquotedblleft{}owner\textquotedblright{}, which reached Marathi through Persian. Its Arabic root is \emph{m-l-k} (Lane). Marathi says \emph{mālak}.
\end{cousinweb}

\begin{grammarlens}[title={What you've built}]
One more word for this chapter.
\end{grammarlens}

\subsection*{Your turn}
\begin{itemize}
\raggedright
  \item \emph{Say it:} \emph{mālak}
  \item \emph{Say it:} \emph{mālak}, once more
  \item \emph{Say it:} say \emph{mālak}
  \item \emph{Recall:} say the Marathi for a right, then the Marathi for a means, then say \emph{mālak} again
\end{itemize}

\begin{tcolorbox}[breakable,colback=teal!4,colframe=teal!35!black,title={Before you move on}]
What does \mr{मालक} mean? (\textquotedblleft{}An owner\textquotedblright{}.) Say it once more.
\end{tcolorbox}
\section[\texorpdfstring{\mr{ताकद} (tākad)}{ताकद (tākad)}]{\mr{ताकद} (tākad) — strength}

\label{lesson:MR-C122-takad}



\begin{quote}
Before the new one: say the Marathi for a means, then the Marathi for an owner.
\end{quote}

\subsection*{You'll want to know: \mr{ताकद}}
\textbf{\mr{ताकद}} — \emph{tākad} — \textquotedblleft{}strength\textquotedblright{}.

\begin{cousinweb}
\textbf{\mr{ताकद}} (\emph{tākad}) is a loan from Arabic \emph{ṭāqa}, \textquotedblleft{}power, strength\textquotedblright{}, which reached Marathi through Persian. Its Arabic root is \emph{ṭ-w-q} (Lane). Persian and Hindustani say \emph{ṭāqat}; Marathi ends it in \textbf{\mr{द}}.
\end{cousinweb}

\begin{grammarlens}[title={What you've built}]
One more word for this chapter.
\end{grammarlens}

\subsection*{Your turn}
\begin{itemize}
\raggedright
  \item \emph{Say it:} \emph{tākad}
  \item \emph{Say it:} \emph{tākad}, once more
  \item \emph{Say it:} say \emph{tākad}
  \item \emph{Recall:} say the Marathi for a means, then the Marathi for an owner, then say \emph{tākad} again
\end{itemize}

\begin{tcolorbox}[breakable,colback=teal!4,colframe=teal!35!black,title={Before you move on}]
What does \mr{ताकद} mean? (\textquotedblleft{}Strength\textquotedblright{}.) Say it once more.
\end{tcolorbox}
\section[\texorpdfstring{\mr{हवा} (havā)}{हवा (havā)}]{\mr{हवा} (havā) — wanted, needed}

\label{lesson:MR-C122-hava}



\begin{quote}
Before the new one: say the Marathi for an owner, then the Marathi for strength.
\end{quote}

\subsection*{You'll want to know: \mr{हवा}}
\textbf{\mr{हवा}} — \emph{havā} — \textquotedblleft{}wanted, needed\textquotedblright{}.

\begin{cousinweb}
The origin of \textbf{\mr{हवा}} (\emph{havā}) is not settled. Neither CDIAL nor DEDR has an entry for it. It is not the Hindi \emph{havā}, \textquotedblleft{}air\textquotedblright{}, which is Arabic.
\end{cousinweb}

\begin{grammarlens}[title={What you've built}]
One more word for this chapter.
\end{grammarlens}

\subsection*{Your turn}
\begin{itemize}
\raggedright
  \item \emph{Say it:} \emph{havā}
  \item \emph{Say it:} \emph{havā}, once more
  \item \emph{Say it:} say \emph{havā}
  \item \emph{Recall:} say the Marathi for an owner, then the Marathi for strength, then say \emph{havā} again
\end{itemize}

\begin{tcolorbox}[breakable,colback=teal!4,colframe=teal!35!black,title={Before you move on}]
What does \mr{हवा} mean? (\textquotedblleft{}Wanted, needed\textquotedblright{}.) Say it once more.
\end{tcolorbox}
\section[\texorpdfstring{\mr{इच्छा} (icchā)}{इच्छा (icchā)}]{\mr{इच्छा} (icchā) — a wish}

\label{lesson:MR-C122-iccha}



\begin{quote}
Before the new one: say the Marathi for strength, then the Marathi for wanted.
\end{quote}

\subsection*{You'll want to know: \mr{इच्छा}}
\textbf{\mr{इच्छा}} — \emph{icchā} — \textquotedblleft{}a wish\textquotedblright{}.

\begin{cousinweb}
\textbf{\mr{इच्छा}} (\emph{icchā}) is Sanskrit \emph{icchā}, \textquotedblleft{}wish\textquotedblright{}, taken over whole as a learned loan, a tatsama (CDIAL 1557). It is built on \emph{iṣ}, \textquotedblleft{}to wish\textquotedblright{}.
\end{cousinweb}

\begin{grammarlens}[title={What you've built}]
One more word for this chapter.
\end{grammarlens}

\subsection*{Your turn}
\begin{itemize}
\raggedright
  \item \emph{Say it:} \emph{icchā}
  \item \emph{Say it:} \emph{icchā}, once more
  \item \emph{Say it:} say \emph{icchā}
  \item \emph{Recall:} say the Marathi for strength, then the Marathi for wanted, then say \emph{icchā} again
\end{itemize}

\begin{tcolorbox}[breakable,colback=teal!4,colframe=teal!35!black,title={Before you move on}]
What does \mr{इच्छा} mean? (\textquotedblleft{}A wish\textquotedblright{}.) Say it once more.
\end{tcolorbox}
